\documentclass[a4paper]{article}
\usepackage{amsfonts}
\usepackage{amsmath}
\usepackage{amssymb}
\usepackage{graphicx}     
\newcommand{\intt}{\displaystyle\int}

\begin{document}
  
\noindent \large Auteur: CouldBeMathijs \\
\noindent \large Studierichting: Informatica\\
\noindent \large Oefeningenreeks 4A nr. 36.4\\

\medskip

\normalsize

$\intt \dfrac{dx}{x^2 + 25} = \dfrac{1}{25} \intt \dfrac{dx}{\left(\dfrac{x}{5}\right)^2 + 1}$\\

Neem $t = \dfrac{x}{5}$ dan $dx = 5dt$\\

$\dfrac{1}{25} \intt \dfrac{5dt}{\left(\dfrac{x}{5}\right)^2 + 1} = \dfrac{1}{5} \intt \dfrac{dt}{\left(\dfrac{x}{5}\right)^2 + 1} = \dfrac{1}{5}\operatorname{Bgtan}\left(\dfrac{x}{5}\right) + c$ met $c \in \mathbb{R}$\\

\begin{thebibliography}{999}
%\bibitem{a} Referentie
\end{thebibliography}
\end{document} 
